\SourceLeaf{046}
\[
\begin{array}{rlllll}
\mathfrak{A}^0A'''&+\mathfrak{B}^0B'''&+\mathfrak{C}^0C'''&+\mathfrak{D}^0&=0\\
&\mathfrak{B}'B'''&+\mathfrak{C}'C'''&+\mathfrak{D}'&=0\\
&&\mathfrak{C}''C'''&+\mathfrak{D}''&=0
\end{array}
\]
unde inveniuntur \(C''',B''',A'''\). Et sic porro.

Uterque modus aeque fere commodus est, si pondera determinationum cunctarum \(x,y,z\) etc. desiderantur; quoties vero e quantitatibus \([\alpha\alpha],\allowbreak[\beta\beta],\allowbreak[\gamma\gamma]\) etc. una tantum vel altera requiritur, manifesto systema prius longe praeferendum erit.

Ceterum combinatio aequationum I cum IV ad easdem formulas perducit, insuperque calculum duplicem ad eruendos valores maxime plausibiles \(A,B,C\) etc. ipsos suppeditat, puta \textit{primo}
\[
\begin{aligned}
A&=-\frac{\mathfrak{L}^0}{\mathfrak{A}^0}-A'\frac{\mathfrak{L}'}{\mathfrak{B}'}-A''\frac{\mathfrak{L}''}{\mathfrak{C}''}-A'''\frac{\mathfrak{L}'''}{\mathfrak{D}'''}-\text{etc.}\\[4pt]
B&=-\frac{\mathfrak{L}'}{\mathfrak{B}'}-B''\frac{\mathfrak{L}''}{\mathfrak{C}''}-B'''\frac{\mathfrak{L}'''}{\mathfrak{D}'''}-\text{etc.}\\[4pt]
C&=-\frac{\mathfrak{L}''}{\mathfrak{C}''}-C'''\frac{\mathfrak{L}'''}{\mathfrak{D}'''}-\text{etc.}\\
&\text{etc.}
\end{aligned}
\]
Calculus alter identicus est cum vulgari, ubi statuitur \(u^0=0,\allowbreak u'=0,\allowbreak u''=0\) etc.

\Article{34}
Quae in art.\ 32 exposuimus, sunt tantummodo casus speciales theorematis generalioris, quod ita se habet:

\textsc{Theorema.} \textit{Designet \(t\) functionem linearem indeterminatarum \(x,y,z\) etc. hanc}
\[
t=fx+gy+hz+\text{etc.}+k
\]
\textit{quae transmutata in functionem indeterminatarum \(u^0,u',u''\) etc. fiat}
\[
t=k^0u^0+k'u'+k''u''+\text{etc.}+K
\]
\textit{Quibus ita se habentibus erit \(K\) valor maxime plausibilis ipsius \(t\), atque pondus huius determinationis}
\[
=\frac1{\mathfrak{A}^0k^0k^0+\mathfrak{B}'k'k'+\mathfrak{C}''k''k''+\text{etc.}}
\]
\textit{Dem.} Pars prior theorematis inde patet, quod valor maxime plausibilis ipsius \(t\) valoribus \(u^0=0,\allowbreak u'=0,\allowbreak u''=0\) etc. respondere debet. Ad posterio\SourceLeaf{047}%
rem demonstrandam observamus, quoniam \(\tfrac12d\Omega=\xi\,dx+\eta\,dy+\zeta\,dz+\text{etc.}\), atque \(dt=f\,dx+g\,dy+h\,dz+\text{etc.}\), esse, pro \(\xi=f.\quad\eta=g,\allowbreak\zeta=h\) etc., independenter a valoribus differentialium \(dx,dy,dz\) etc.
\[
d\Omega=2dt
\]
Hinc vero sequitur, pro iisdem valoribus \(\xi=f,\allowbreak\eta=g,\allowbreak\zeta=h\) etc., fieri
\[
\begin{aligned}
&\frac{u^0}{\mathfrak{A}^0}du^0+\frac{u'}{\mathfrak{B}'}du'+\frac{u''}{\mathfrak{C}''}.du''+\text{etc.}\\
&\qquad=k^0du^0+k'du'+k''du''+\text{etc.}
\end{aligned}
\]
Iam facile perspicitur, si \(dx,dy,dz\) etc. sint ab invicem independentes, etiam \(du^0,du',du''\) etc., ab invicem independentes esse; unde colligimus, pro \(\xi=f,\allowbreak\eta=g,\allowbreak\zeta=h\) etc. esse
\[
u^0=\mathfrak{A}^0k^0,\quad u'=\mathfrak{B}'k',\quad u''=\mathfrak{C}''k''\quad\text{etc.}
\]
Quamobrem valor ipsius \(\Omega\), iisdem valoribus respondens erit
\[
=\mathfrak{A}^0k^0k^0+\mathfrak{B}'k'k'+\mathfrak{C}''k''k''+\text{etc.}+M
\]
unde per art.\ 29 theorematis nostri veritas protinus demanat.

Ceterum si transformationem functionis \(t\) immediate, i.\ e. absque cognitione substitutionum IV. art.\ 32, perficere cupimus, praesto sunt formulae:
\[
\begin{aligned}
f&=\mathfrak{A}^0k^0\\
g&=\mathfrak{B}^0k^0+\mathfrak{B}'k'\\
h&=\mathfrak{C}^0k^0+\mathfrak{C}'k'+\mathfrak{C}''k''\quad\text{etc.,}
\end{aligned}
\]
unde coëfficientes \(k^0,k',k''\) etc. deinceps determinabuntur, tandemque habebitur
\[
K=k-\mathfrak{L}^0k^0-\mathfrak{L}'k'-\mathfrak{L}''k''-\text{etc.}
\]

\Article{35}
Tractatione peculiari dignum est problema sequens, tum propter utilitatem practicam, tum propter solutionis concinnitatem.

\textit{Invenire mutationes valorum maxime plausibilium incognitarum ab accessione aequationis novae productas, nec non pondera novarum determinationum.}

Retinebimus designationes in praecedentibus adhibitas, ita ut aequationes primitivae, ad pondus \(=1\) reductae, sint hae \(v=0,\allowbreak v'=0,\allowbreak v''=0\) etc.; ag\SourceLeaf{048}%
gregatum indefinitum \(vv+v'v'+v''v''\quad\text{etc.}=\Omega\); porro ut \(\xi,\eta,\zeta\) etc. sint quotientes differentiales partiales
\[
\frac{d\Omega}{2dx},\quad\frac{d\Omega}{2dy},\quad\frac{d\Omega}{2dz}\quad\text{etc.}
\]
denique ut ex eliminatione indefinita sequatur
\[
\left.\begin{aligned}
x&=A+[\alpha\alpha]\xi+[\alpha\beta]\eta+[\alpha\gamma]\zeta+\text{etc.}\\
y&=B+[\alpha\beta]\xi+[\beta\beta]\eta+[\beta\gamma]\zeta+\text{etc.}\\
z&=C+[\alpha\gamma]\xi+[\beta\gamma]\eta+[\gamma\gamma]\zeta+\text{etc.}
\end{aligned}\right\}\ (\mathrm{I})
\]
Iam supponamus, accedere aequationem novam \(v^*=0\) (proxime veram, et cuius pondus \(=1\)), et inquiramus, quantas mutationes hinc nacturi sint tum valores incognitarum maxime plausibiles \(A,B,C\) etc., tum coëfficientes \([\alpha\alpha],[\alpha\beta]\) etc.

Statuamus \(\Omega+v^*v^*=\Omega^*,\allowbreak\frac{d\Omega^*}{2dx}=\xi^*,\allowbreak\frac{d\Omega^*}{2dy}=\eta^*,\allowbreak\frac{d\Omega^*}{2dz}=\zeta^*\) etc. supponamusque, hinc per eliminationem sequi
\[
x=A^*+[\alpha\alpha^*]\xi^*+[\alpha\beta^*]\eta^*+[\alpha\gamma^*]\zeta^*\quad\text{etc.}
\]
Denique sit
\[
v^*=fx+gy+hz+\text{etc.}+k
\]
prodeat inde, substitutis pro \(x,y,z\) etc. valoribus ex (I),
\[
v^*=F\xi+G\eta+H\zeta+\text{etc.}+K
\]
statuaturque
\[
Ff+Gg+Hh+\text{etc.}=\omega
\]
Manifesto \(K\) erit valor maxime plausibilis functionis \(v^*\), quatenus ex aequationibus primitivis sequitur, sine respectu valoris \(0\), quem observatio accessoria praebuit, atque \(\frac1\omega\) pondus istius determinationis.

Iam habemus
\[
\xi^*=\xi+fv^*,\quad\eta^*=\eta+gv^*,\quad\zeta^*=\zeta+hv^*\quad\text{etc.}
\]
adeoque
\[
\begin{aligned}
&F\xi^*+G\eta^*+H\zeta^*+\text{etc.}+K\\
&\qquad=v^*(1+Ff+Gg+Hh+\text{etc.})
\end{aligned}
\]
sive
\[
v^*=\frac{F\xi^*+G\eta^*+H\zeta^*+\text{etc.}+K}{1+\omega}
\]
